\ExplSyntaxOn
\NewExpandableDocumentCommand{\Lang}{m m}
  { \bool_if:NTF \g_script_english_bool {#2} {#1} }
\NewExpandableDocumentCommand{\IfClozeTF}{m m}
  { \bool_if:NTF \g_script_cloze_bool {#1} {#2} }
\NewExpandableDocumentCommand{\IfEnglishTF}{m m}
  { \bool_if:NTF \g_script_english_bool {#1} {#2} }
\ExplSyntaxOff

\NewDocumentCommand{\EnglishTerm}{m}{%
  {\color{black!55}\itshape(engl.: #1)}%
}
\NewDocumentCommand{\GermanTerm}{m}{%
  {\color{black!55}\itshape(German: #1)}%
}

% Use a prebuilt standalone PDF when available; direct main.tex builds fall
% back to the TikZ source and therefore remain self-contained.
\makeatletter
\NewDocumentCommand{\CachedTikzFigure}{m}{%
  \begingroup
    \filename@parse{#1}%
    \edef\CachedTikzPdf{\TikzCacheRoot/\filename@base.pdf}%
    \IfFileExists{\CachedTikzPdf}
      {\includegraphics{\CachedTikzPdf}}
      {\input{#1}}%
  \endgroup
}
\makeatother

% Right-hand operation in an equation transformation chain.
\NewDocumentCommand{\TransformOp}{m}{%
  {\color{black!55}\Big|\;#1}%
}

% Uniform overview of the three possible solution types of a linear system.
% The starred form crosses out the impossible "no solution" case and is used
% for homogeneous systems.
\NewDocumentCommand{\SolutionTypeBox}{m m}{%
  \begin{tikzpicture}[baseline=(solutiontype.base)]
    \node[
      draw=black,
      line width=0.4pt,
      text width=0.24\linewidth,
      inner xsep=0pt,
      inner ysep=0.3em,
      minimum height=2.1em,
      align=center
    ] (solutiontype) {#2};
    \IfBooleanT{#1}{%
      \draw[red!75!black,line width=1.6pt,shorten <=-5pt,shorten >=-5pt]
        (solutiontype.south west) -- (solutiontype.north east);
      \draw[red!75!black,line width=1.6pt,shorten <=-5pt,shorten >=-5pt]
        (solutiontype.north west) -- (solutiontype.south east);
    }%
  \end{tikzpicture}%
}
\NewDocumentCommand{\SolutionTypeCases}{s}{%
  \IfBooleanTF{#1}{\par\medskip}{\par\smallskip}%
  \noindent
  \SolutionTypeBox{\BooleanFalse}{\Lang{genau eine}{exactly one}}%
  \hfill
  \SolutionTypeBox{#1}{\Lang{keine}{none}}%
  \hfill
  \SolutionTypeBox{\BooleanFalse}{\Lang{$\infty$-viele}{$\infty$ many}}%
  \IfBooleanTF{#1}{\par\medskip}{\par\smallskip}%
}

% Restrained CD background for short answers, kept light for pencil writing.
\definecolor{ScriptClozeBackground}{HTML}{E5E4D8}
\NewDocumentCommand{\Cloze}{m m}{%
  \IfClozeTF
    {{\setlength{\fboxsep}{0pt}%
      \colorbox{ScriptClozeBackground}{\makebox[#1][c]{\strut}}}}
    {#2}%
}

\NewDocumentCommand{\FrontMatterChapter}{m}{%
  \chapter*{#1}%
  \addcontentsline{toc}{chapter}{#1}%
}

% Grid dimensions are specified as numbers of 5 mm squares:
% \begin{ClozeArea}[<columns>]{<rows>} ... \end{ClozeArea}
\NewDocumentEnvironment{ClozeArea}{O{34} m +b}{%
  \par\medskip\noindent
  \pgfmathsetlengthmacro{\ClozeAreaWidth}{#1*5mm}%
  \pgfmathsetlengthmacro{\ClozeAreaHeight}{#2*5mm}%
  \hfill
  \IfClozeTF{%
    \begin{tikzpicture}
      \path[use as bounding box]
        (0,0) rectangle (\ClozeAreaWidth,\ClozeAreaHeight);
      \draw[step=5mm,black!14,line width=0.2pt]
        (0,0) grid (\ClozeAreaWidth,\ClozeAreaHeight);
      \draw[black!25,line width=0.3pt]
        (0,0) rectangle (\ClozeAreaWidth,\ClozeAreaHeight);
    \end{tikzpicture}%
  }{%
    \begin{minipage}{\ClozeAreaWidth}#3\end{minipage}%
  }%
  \hfill\mbox{}%
  \par\medskip
}{}

% Grid that continues a visible equation within the same display.
% Dimensions again count exact 5 mm squares.
\NewDocumentEnvironment{ClozeContinuation}{O{12} m +b}{%
  \pgfmathsetlengthmacro{\InlineClozeAreaWidth}{#1*5mm}%
  \pgfmathsetlengthmacro{\InlineClozeAreaHeight}{#2*5mm}%
  \IfClozeTF{%
    \vcenter{\hbox{%
      \begin{tikzpicture}
        \path[use as bounding box]
          (0,0) rectangle (\InlineClozeAreaWidth,\InlineClozeAreaHeight);
        \draw[step=5mm,black!14,line width=0.2pt]
          (0,0) grid (\InlineClozeAreaWidth,\InlineClozeAreaHeight);
        \draw[black!25,line width=0.3pt]
          (0,0) rectangle (\InlineClozeAreaWidth,\InlineClozeAreaHeight);
      \end{tikzpicture}%
    }}%
  }{#3}%
}{}

% Shared geometry for branching and merging equations.
\newlength{\EquationTemplateUpperWidth}
\newlength{\EquationTemplateLowerWidth}
\newlength{\EquationTemplateMaxWidth}
\newlength{\EquationTemplateUpperNoteWidth}
\newlength{\EquationTemplateLowerNoteWidth}
\newlength{\EquationTemplateMaxNoteWidth}
\NewDocumentCommand{\EquationTemplateMeasure}{m m}{%
  \settowidth{\EquationTemplateUpperWidth}{\ensuremath{\displaystyle #1}}%
  \settowidth{\EquationTemplateLowerWidth}{\ensuremath{\displaystyle #2}}%
  \ifdim\EquationTemplateUpperWidth>\EquationTemplateLowerWidth
    \setlength{\EquationTemplateMaxWidth}{\EquationTemplateUpperWidth}%
  \else
    \setlength{\EquationTemplateMaxWidth}{\EquationTemplateLowerWidth}%
  \fi
}
\NewDocumentCommand{\EquationTemplateMeasureNotes}{m m}{%
  \settowidth{\EquationTemplateUpperNoteWidth}{#1}%
  \settowidth{\EquationTemplateLowerNoteWidth}{#2}%
  \ifdim\EquationTemplateUpperNoteWidth>\EquationTemplateLowerNoteWidth
    \setlength{\EquationTemplateMaxNoteWidth}{\EquationTemplateUpperNoteWidth}%
  \else
    \setlength{\EquationTemplateMaxNoteWidth}{\EquationTemplateLowerNoteWidth}%
  \fi
}
\NewDocumentCommand{\EquationTemplateOverlay}{m m m}{%
  \makebox[0pt][l]{%
    \hspace{#2}\raisebox{#1}[0pt][0pt]{#3}}%
}
\NewDocumentCommand{\EquationTemplateRelation}{m m}{%
  \rotatebox{#2}{\ensuremath{\displaystyle #1}}%
}
\NewDocumentCommand{\EquationTemplateEquals}{m}{%
  \EquationTemplateRelation{=}{#1}%
}
\NewDocumentCommand{\EquationTemplateUse}{m}{%
  \ifmmode
    \vcenter{\hbox{#1}}%
  \else
    #1%
  \fi
}
\NewDocumentCommand{\MathHighlight}{m}{%
  \begingroup
  \setlength{\fboxsep}{1.5pt}%
  \colorbox{yellow!65}{\ensuremath{\displaystyle #1}}%
  \endgroup
}

% One expression continued in two alternative directions:
% \EquationBranches[<relation>]{<start>}{<upper continuation>}{<upper note>}
%                              {<lower continuation>}{<lower note>}
\NewDocumentCommand{\EquationBranchesContent}{O{=} m m m m m}{%
  \begingroup
  \EquationTemplateMeasure{#3}{#5}%
  \EquationTemplateMeasureNotes{#4}{#6}%
  \ensuremath{\displaystyle #2}%
  \hspace{0.15cm}%
  \EquationTemplateOverlay{0.7em}{0pt}{\EquationTemplateRelation{#1}{35}}%
  \EquationTemplateOverlay{1.4em}{0.55cm}{%
    \makebox[\EquationTemplateMaxWidth][l]{\ensuremath{\displaystyle #3}}%
    \hspace{0.55cm}#4}%
  \EquationTemplateOverlay{-0.7em}{0pt}{\EquationTemplateRelation{#1}{-35}}%
  \EquationTemplateOverlay{-1.4em}{0.55cm}{%
    \makebox[\EquationTemplateMaxWidth][l]{\ensuremath{\displaystyle #5}}%
    \hspace{0.55cm}#6}%
  \rule[-1.8em]{0pt}{4.3em}%
  % Reserve the full visible branch width so the complete construction,
  % including both notes, is centered by the surrounding display.
  \hspace{1.10cm}%
  \hspace{\EquationTemplateMaxWidth}%
  \hspace{\EquationTemplateMaxNoteWidth}%
  \endgroup
}

\NewDocumentCommand{\EquationBranches}{O{=} m m m m m}{%
  \EquationTemplateUse{%
    \EquationBranchesContent[#1]{#2}{#3}{#4}{#5}{#6}}%
}

% Two alternative expressions leading to one common result:
% \EquationMerge{<upper note>}{<upper expression>}{<lower note>}
%                {<lower expression>}{<result>}
\NewDocumentCommand{\EquationMergeContent}{m m m m m}{%
  \begingroup
  \EquationTemplateMeasure{#2}{#4}%
  \makebox[\EquationTemplateMaxWidth][l]{%
    \EquationTemplateOverlay{1.4em}{0pt}{%
      \makebox[0pt][r]{#1\hspace{0.55cm}}%
      \makebox[\EquationTemplateMaxWidth][r]{\ensuremath{\displaystyle #2}}}%
    \EquationTemplateOverlay{-1.4em}{0pt}{%
      \makebox[0pt][r]{#3\hspace{0.55cm}}%
      \makebox[\EquationTemplateMaxWidth][r]{\ensuremath{\displaystyle #4}}}%
    \rule[-1.8em]{0pt}{4.3em}}%
  \hspace{0.15cm}%
  \EquationTemplateOverlay{0.7em}{0pt}{\EquationTemplateEquals{-35}}%
  \EquationTemplateOverlay{-0.7em}{0pt}{\EquationTemplateEquals{35}}%
  \hspace{0.55cm}%
  \ensuremath{\displaystyle #5}%
  \endgroup
}

\NewDocumentCommand{\EquationMerge}{m m m m m}{%
  \EquationTemplateUse{%
    \EquationMergeContent{#1}{#2}{#3}{#4}{#5}}%
}

\newcommand{\R}{\mathbb{R}}
\newcommand{\N}{\mathbb{N}}
\newcommand{\Q}{\mathbb{Q}}
\newcommand{\Z}{\mathbb{Z}}
\newcommand{\C}{\mathbb{C}}
% Right-angle mark: a circular arc and a small dot inside its sector.
% The angles describe the visible projection; a spatial right angle need not
% appear as 90 degrees on the page. Radius and style remain figure-specific.
% Arguments: [style]{vertex}{start angle}{end angle}{arc radius}
%            [dot radius][dot distance as a fraction of the arc radius].
\NewDocumentCommand{\RightAngleMark}{O{black!65} m m m m O{0.85pt} O{0.6667}}{%
  \pgfmathsetmacro{\RightAngleMidAngle}{(#3+#4)/2}%
  \pgfmathsetmacro{\RightAngleDotDistance}{#5*(#7)}%
  \draw[#1,line width=0.75pt] ($#2+(#3:#5)$)
    arc[start angle=#3,end angle=#4,radius=#5];%
  \fill[#1] ($#2+(\RightAngleMidAngle:\RightAngleDotDistance)$) circle (#6);%
}
\DeclareMathOperator{\grad}{grad}
\DeclareMathOperator{\arsinh}{arsinh}
\DeclareMathOperator{\arcosh}{arcosh}
\DeclareMathOperator{\artanh}{artanh}
\newcommand{\corresponds}{\mathrel{\widehat{=}}}
